% Template authors: Bogumił Kamiński and Michał Jakubczyk

\documentclass[12pt,a4paper,twoside,openany]{book}
\usepackage[T1]{fontenc}
\usepackage[utf8]{inputenc}
\usepackage[polish]{babel}
\usepackage{graphicx}
\usepackage{times}
\usepackage{indentfirst}
\usepackage[left=3cm,right=2cm,top=2.5cm,bottom=2.5cm]{geometry}
\usepackage{natbib}
\usepackage{enumitem}
\usepackage{color}
\usepackage{soul}
\usepackage{tikz}
\usepackage{url}
\usepackage{todonotes}
\usepackage{verbatim}
\setlist{itemsep=0pt}
\setlist{nolistsep}
\frenchspacing
\linespread{1.3}
\addto\captionspolish{%
\renewcommand*\listtablename{Spis tabel}
\renewcommand*\tablename{Tabela}
}
\usepackage{titlesec}
\titlelabel{\thetitle.\quad}

\frenchspacing

\begin{document}

\begin{center}
\includegraphics[scale=0.3]{sgh_full.png}

\vspace{1cm}

% tu i dalej fbox należy usunąć i wpisać odpowiednią wartość
Studium magisterskie
\end{center}

\vspace{1cm}

\noindent Kierunek: Analiza danych - big data

\noindent Specjalność: \fbox{\vphantom{I}\dots}

\vspace{1cm}

{
\leftskip=10cm\noindent
Aleksander Wojnarowicz\newline
Nr albumu: 77438

}

\vspace{2cm}

\title{Wzór pracy dyplomowej SGH}
\makeatletter

\begin{center}
\LARGE\bf
\fbox{\@title}
\end{center}

\vspace{2cm}

{
\leftskip=10cm\noindent
Praca magisterska
napisana w~Instytucie Ekonometrii\newline
pod kierunkiem naukowym\newline
\fbox{promotor}

}

\vfill

\begin{center}
Warszawa, \the\year
\end{center}
\thispagestyle{empty}

\clearpage
\thispagestyle{empty}
\mbox{}
% druga strona będzie pusta, ponieważ drukujemy dwustronnie
% a mbox jest po to, żeby ta strona się pokazała
% od procenta robimy komentarze
\clearpage

\tableofcontents

\clearpage

\chapter*{Uwagi techniczne, skasuj rozdział po uwzględnieniu}

Wymagane jest wgranie pracy do systemu Overleaf (\url{https://www.overleaf.com/}) i udostępnienie promotorowi z~prawem zmian.

Pytania do promotora w~tekście proszę zadawać \todo[inline]{o tak}, tj.~stosując \verb!\todo[inline]{o tak}!.

Po kropkach niekończących zdania (np.~po `np.', `tj.', `itd.', itd.) stawiamy znak tyldy, tj.~\verb!~! (żeby zmniejszyć odstęp). Wykorzystujemy także tyldę między odwołaniem i~numerem, np.~`Tabela~3' (oczywiście odwołanie robione automatycznie przez \verb!\ref{}!). I wreszcie, stosujemy tyldę po pojedynczych literach (np.~`i', `a'), bo to niełamliwa spacja i~gwarantuje brak takich liter na końcu linii.

Przykłady odwołań źrółowych:
\begin{itemize}
\item \citep{dolan2000}, \citet{dolan2000}, \defcitealias{dolan2000}{Dolana (2000)}\citetalias{dolan2000};
\item \citep{wakker1999}, \citet{wakker1999}, \defcitealias{wakker1999}{Wakkera i Zhanka (2000)}\citetalias{wakker1999};
\item \citep{drummond1997}, \citet{drummond1997}, \defcitealias{drummond1997}{Drummonda et al.~(2000)}\citetalias{drummond1997}.
\end{itemize}

Co warto przeczytać o~\LaTeX:
\begin{enumerate}
\footnotesize % tak zmieniamy rozmiar fontu
\item \url{http://www.tex.ac.uk/ctan/info/gentle/gentle.pdf},
\item \url{ftp://sunsite.icm.edu.pl/pub/CTAN/info/lshort/english/lshort.pdf},
\item \url{http://paws.wcu.edu/tsfoguel/tikzpgfmanual.pdf},
\item \url{https://www.overleaf.com/latex/learn/free-online-introduction-to-latex-part-1}.
\end{enumerate}

\chapter{Wprowadzenie}

Nierelatywistyczna fizyka klasyczna stanowi podstawę obrazu świata, postrzeganego oczami ludzkiej wyobraźni. Umysł homo sapiens wyewoluował w otoczeniu obiektów makroskopowych, a ponadto nigdy nie doświadczył (ani nawet nie był w stanie sobie wyobrazić) prędkości choć trochę zbliżonych do prędkości światła. Z tychże powodów, dwa wielkie osiągnięcia współczesnej fizyki – mechanika kwantowa oraz teoria względności (zarówno ogólna, jak i szczególna) – są tak nieintuicyjne. Tymczasem płynące z nich wnioski i wyniki stanowią podstawę dzisiejszej cywilizacji. Dylatacja czasu jest istotnym czynnikiem, jaki muszą uwzględniać precyzyjne systemy nawigacyjne GPS. Tranzystor z kolei – podstawowy budulec układów scalonych, a zatem również wszechobecnej dziś elektroniki – powstał i był doskonalony dzięki badaniom nad kwantowymi własnościami materii w nanoskali.

\section{Cel pracy}
Celem niniejszej pracy jest zbadanie oraz przeanalizowanie pewnego interesującego zastosowania fizyki mikroświata w dziedzinie finansów, a konkretnie - budowie portfela inwestycyjnego. Poza oceną samych rezultatów opisanej w kolejnych rozdziałach procedury istotne jest też porównanie ich do wyników jakie można otrzymać, stosując rozwiązania klasyczne tj. zaimplementowane na komputerze niekwantowym, który ponadto nie symuluje komputera kwantowego.

Quantum Computing (QC) ...
\section{Granice klasycznej informatyki}

\subsection{Tranzystory i elektronika}
Podstawą działania nie tylko komputerów, ale w zasadzie niemal całej współczesnej elektroniki są tranzystory – mikroskopijne „przełączniki”, które, w zależności od stanu w jakim się znajdują, mogą albo przepuszczać, albo zatrzymywać elektrony płynące przez obwody elektryczne. Te dwa stany koduje się jako 0 (przełącznik nie przepuszczający elektronów) oraz 1 (przełącznik przepuszczający elektrony). Takie podejście stanowi bazę kodu binarnego, na którym oparte jest działanie komputerów.
Na bazie tranzystorów konstruowane są bramki logiczne, które z kolei pozwalają na przeprowadzanie podstawowych operacji arytmetycznych, wykonywanych na wyżej wspomnianym kodzie binarnym. Poprzez składanie ze sobą wielu różnych tego typu operacji otrzymujemy pożądane zachowania maszyn elektronicznych – obliczenia naukowe, grafikę komputerową, komunikację internetową etc.
	Celem zwiększenia mocy obliczeniowych maszyny elektronicznej, konieczne jest zwiększenie liczby tranzystorów, które zawiera w sobie ta maszyna. Istotne jest jednak, aby nie wpłynęło to zbytnio na rozmiar urządzenia – powinno być najmocniejsze jak to tylko możliwe, będąc przy tym najmniejsze jak to tylko możliwe. Z tego względu na przestrzeni lat następował konsekwentny spadek rozmiaru tranzystora (patrz. Tabela 1).
	


\begin{table}[h]
\centering
\caption{Tranzystory na przestrzeni lat}
\label{tab:przyklad}
\footnotesize
\begin{tabular}{|l|r|}
\hline
Rok & Rozmiar\\
\hline
1971 & 10000  \\
1974 & 6000  \\
1977 & 3000 \\
1981 & 1500 \\
1984 & 1000 \\
1987 & 800 \\
1990 & 600 \\
1993 & 350 \\
1996 & 250 \\
1999 & 180 \\
2001 & 130 \\
2003 & 90 \\
2005 & 65 \\
2007 & 45 \\
2009 & 32 \\
2012 & 22 \\
2014 & 14 \\
2016 & 10 \\
2018 & 7 \\
2020 & 5 \\

\hline
\end{tabular} 
\end{table}

\subsection{Tunelowanie kwantowe}

Ten postępujący wzrost mocy obliczeniowej ma jednak granicę, wyznaczoną przez prawa fizyki kwantowej. W miarę jak tranzystory kurczą się coraz bardziej do rozmiarów porównywalnych z rozmiarami atomów, coraz istotniejszą rolę gra zjawisko tunelowania kwantowego. Z pewnym prawdopodobieństwem, które rośnie odwrotnie do rozmiarów „przełącznika”, elektron jest w stanie pokonać barierę potencjału, nawet przy braku odpowiedniej ilości energii, której wymaga do tego fizyka klasyczna. Innymi słowy, istnieje zawsze niezerowe prawdopodobieństwo, że pomimo iż tranzystor będzie w stanie 0, to elektrony i tak będą mogły się przez niego przedostać . Oczywistą konsekwencją tego zjawiska jest niepoprawne działanie elektroniki – wyniki pracy maszyn nabrałyby charakteru probabilistycznego, a zatem nie można by było na nich polegać.
	
Prawdopodobieństwo zajścia kwantowego tunelowania zadane jest, przy pewnych założeniach , przez wzór :
\begin{equation}
e^{-\frac{4a\pi}{h}\sqrt{2m(V - E)}}
\end{equation}
gdzie: \\
\indent a – grubość bariery (w naszym przypadku – rozmiar tranzystora) \\
\indent h – stała Planka ($ \approx 6.626 \cdots 10^{-34}  kg  \cdot m^2 / s$) \\
\indent m – masa cząstki (dla elektronu $m \approx 9,109 3837015 \cdot 10^{-31} kg$) \\
\indent V – energia potencjalna bariery \\ 
\indent E – energia kinetyczna cząsteczki ( dla zjawiska tunelowania kwantowego E < V) \\

Podstawiając za wartość energii elektronu 4.5 eV, a za wartość energii potencjalnej bariery 5 eV otrzymamy prawdopodobieństwo przejścia elektronu przez barierę o rozmiarze 5nm wynoszące około 2 x 10-16. To bardzo mała wartość, jednak jeśli zmniejszymy rozmiar bariery do 0.5 nm (500 pm) – prawdopodobieństwo to przekroczy już 0.25 – wartość zdecydowanie nieakceptowalna z punktu widzenia praktycznych zastosowań. Jest więc wyraźnie widoczne, że już w niedalekiej przyszłości przestanie być możliwe zwiększanie mocy obliczeniowej komputerów w sposób, w jaki odbywa się to obecnie. Pojawia się zatem potrzeba zbadania alternatywnych metod, które pozwoliłyby ulepszać pracę urządzeń elektrycznych, bez nadwątlania rzetelności zwracanych przez nie rezultatów.



\chapter{Kolejny rozdział}

Wstępik, tj.~co jest w tym rozdziale.

\section{Podrozdział}



\chapter{I następny}
\label{sec:nast}

Bieżący rozdział (czyli rozdział~\ref{sec:nast}) zawiera odwołanie do samego siebie. Stosujmy odwołania automatyczne.

Dodatkowo w niniejszym rozdziale zamieszczono przykładową tabelę (uwaga --- \LaTeX\ umiejscawia ją dynamicznie).

% warto w kodzie poniżej zwrócić uwagę na prawidłowy sposób formatowania liczb z przecinkiem jako znakiem dziesiętnym
% a dodatkowo na wyrównywanie liczb ułatwiające porównywanie


\begin{table}[h]
\centering
\caption{Przykładowa tabela.}
\label{tab:przyklad}
\footnotesize
\begin{tabular}{|l|r|r|}
\hline
Wariant & $N=5$ & $N=10$\\
\hline
A & $2$ & $3\phantom{{,}1}$ \\
B & $5$ & $3{,}1$ \\

\hline
\end{tabular} 
\end{table}

\clearpage

\chapter{Matematyka}

Wzory proste umieszczamy w tekście $a^2+b^2=c^2$, wzory bardziej skomplikowane --- poza nim:
\begin{equation}
\sum_{n=1}^{+\infty}\frac{1}{n^2}=\frac{\pi^2}{6}.
\end{equation}

\section{Pierwszy podrozdział}

\begin{figure}[h]
  \centering
  \begin{tikzpicture}[scale=4]
    \draw[->] (0,0) -- (1.4,0) node[anchor=west] {$n$};
    \draw[->] (0,0) -- (0,1.35) node[anchor=south] {Iloraz};
    \foreach \y in {0, 0.25,0.5,0.75,1, 1.25}
      \draw (0.03, \y) -- (-0.03, \y) node[anchor=east] {$\y$};
     \foreach \x in {0, 2, 4, 6, 8, 10, 12}
       \draw ({\x/10}, 0.03) -- ({\x/10}, -0.03) node[anchor=north] {$\x$};

       \draw (0.200000,1.164983) circle[radius=0.02];
       \draw (0.300000,1.176768) circle[radius=0.02];
       \draw (0.400000,0.877412) circle[radius=0.02];
       \draw (0.500000,0.456929) circle[radius=0.02];
       \draw (0.600000,0.351427) circle[radius=0.02];
       \draw (0.700000,0.294229) circle[radius=0.02];
       \draw (0.800000,0.246016) circle[radius=0.02];
       \draw (0.900000,0.224231) circle[radius=0.02];
       \draw (1.000000,0.168351) circle[radius=0.02];
       \draw (1.100000,0.191722) circle[radius=0.02];
       \draw (1.200000,0.173029) circle[radius=0.02];
       \draw (1.300000,0.176309) circle[radius=0.02];

       \fill (0.200000,1.225646) circle[radius=0.02];
       \fill (0.300000,0.986035) circle[radius=0.02];
       \fill (0.400000,0.945566) circle[radius=0.02];
       \fill (0.500000,0.688660) circle[radius=0.02];
       \fill (0.600000,0.556635) circle[radius=0.02];
       \fill (0.700000,0.483555) circle[radius=0.02];
       \fill (0.800000,0.443744) circle[radius=0.02];
       \fill (0.900000,0.413361) circle[radius=0.02];
       \fill (1.000000,0.344671) circle[radius=0.02];
       \fill (1.100000,0.374543) circle[radius=0.02];
       \fill (1.200000,0.358811) circle[radius=0.02];
       \fill (1.300000,0.362736) circle[radius=0.02];
  \end{tikzpicture}
  \caption{Oto przykładowy, skomplikowany rysunek. Często będzie łatwiej. \label{fig:algcomparison}}
\end{figure}

Każdy rysunek przywołujemy w tekście, także rysunek \ref{fig:algcomparison}. Można też ładować obrazy, jak pokazano na stronie tytułowej.

\section{Drugi}

\section{Trzeci}

\clearpage

\chapter{Podsumowanie}

\clearpage
\addcontentsline{toc}{chapter}{Bibliografia}
\begin{thebibliography}{99}
\setlength{\itemsep}{0pt}%
\bibitem[Dolan(2000)]{dolan2000} Dolan P.~(2000), The Measurement of (...), w: A.J.~Culyer, J.P.~Newhouse (ed.), Handbook of Health Economics, Volume 1B, Elsevier, s.~1723--1760
\bibitem[Wakker i Zhank(1999)]{wakker1999} Wakker P., H.~Zank (1999), A Unified Derivation (...), \textit{Journal of  Mathematical Economics}, 32, s.~1-19
\bibitem[Drummond et al.(1997)]{drummond1997} Drummond M.F., B.J.~O'Brien, G.L.~Stoddart, G.W.~Torrance (1997), Methods for Economic Evaluation (...), Oxford University Press
\end{thebibliography}

\clearpage
\addcontentsline{toc}{chapter}{Spis rysunków}
\listoffigures

\clearpage
\listoftables
\addcontentsline{toc}{chapter}{Spis tabel}

\appendix
\chapter*{Kody źródłowe}
\addcontentsline{toc}{chapter}{Kody źródłowe}

\section*{Analiza 1}
\begin{verbatim}
x <- 1:10
y <- x + rnorm(10)
summary(lm(y ~ x))
\end{verbatim}

\section*{Analiza 2}
\begin{verbatim}
x <- rnorm(10000)
plot(density(x))
\end{verbatim}

\clearpage

\chapter*{Streszczenie}
\addcontentsline{toc}{chapter}{Streszczenie}

Tutaj piszemy streszczenie. Między 200 a 350 słów. Nie jest omówieniem struktury pracy. Zawiera cel, metodę, dane, wyniki, wnioski. Nie ma odwołań źródłowych, list, wykresów, tabel. Ma być zrozumiałe dla osoby nieczytającej pracy.

\end{document}

